% ==================== 正文章节编号、标题间距及章节页眉页脚（含 plain 样式） ====================

%%%%%%%% 章节 %%%%%%%%

% 为 subsubsection 也启用编号
\setcounter{secnumdepth}{3}

% 章节标题字体
\newlength{\xmu@chaptertitleheight}
% 章号与章名保持同一行；标题区域至少占两个正文行高，长标题自然增高。
\newcommand{\xmu@chaptertitlebox}[1]{%
    \setbox\@tempboxa=\vbox{\hsize=\linewidth\centering
        \ifttl@label
            \ifcsdef{xmu@arabicchapter}{\thechapter}{第\zhnumber{\thechapter}章}\hspace{1em}%
        \fi
        #1\par
    }%
    \ifdim\dimexpr\ht\@tempboxa+\dp\@tempboxa\relax<\xmu@chaptertitleheight
        \vbox to\xmu@chaptertitleheight{\vfil\box\@tempboxa\vfil}%
    \else
        \box\@tempboxa
    \fi
}
\titleformat{\chapter}[block]{\xmu@heiti\zihao{-3}\linespread{1.5}\selectfont\let\ttl@strut\@empty}{}{0pt}{\xmu@chaptertitlebox}
\titleformat{\section}{\xmu@heiti\zihao{4}}{\thesection}{0.5em}{}
\titleformat{\subsection}{\xmu@heiti\zihao{-4}}{\thesubsection}{0.5em}{}
\titleformat{\subsubsection}{\sf\bfseries\zihao{-4}}{\thesubsubsection}{0.5em}{}

% 章节标题行间距
\titlespacing{\chapter}{\z@}{0pt}{0pt}
\titlespacing{\section}{\z@}{18bp}{12bp}
\titlespacing{\subsection}{\z@}{12bp}{12bp}
\titlespacing{\subsubsection}{\z@}{1ex \@plus 0.7ex \@minus 0ex}{0.7ex}

% 章节页眉、页脚
% 普通段落统一设置，标题、目录、图表及参考文献在各自模块中单独处理。
\newcommand{\xmu@bodyparagraphs}{%
    \xmu@songti\zihao{-4}%
    \linespread{1.5}\selectfont
    \setlength{\xmu@chaptertitleheight}{2\baselineskip}%
    \setlength{\parindent}{2em}%
    \setlength{\parskip}{0pt}%
    \setlength{\leftskip}{0pt}%
    \setlength{\rightskip}{0pt}%
    \setlength{\parfillskip}{0pt plus 1fil}%
    \setlength{\mathindent}{2em}%
}
\AtBeginDocument{\xmu@bodyparagraphs}

% 打印版按需补空白页，补页不显示页眉、页脚和页码。
\renewcommand{\cleardoublepage}{%
    \clearpage
    % 封面后或正文前会重置逻辑页码，打印版应依据实际纸张页序判断奇偶。
    \ifcsdef{c@abspage}{\@tempcnta=\value{abspage}}{\@tempcnta=\c@page}%
    \if@twoside\ifodd\@tempcnta\else
        \hbox{}\thispagestyle{empty}\newpage
        \if@twocolumn\hbox{}\thispagestyle{empty}\newpage\fi
    \fi\fi
}
\fancypagestyle{xmu@front}{%
    \fancyhf{}%
    \renewcommand{\headrulewidth}{0.4pt}%
    \fancyhead[OC]{\xmu@songti\zihao{-5}\leftmark}%
    \fancyhead[EC]{\xmu@songti\zihao{-5}\xmu@title}%
    \fancyfoot[C]{\normalfont\rmfamily\zihao{-5}\thepage}%
}
% 摘要和目录复用固定 36 磅标题行距、24 磅段后间距，字体由调用处保留。
\newcommand{\xmu@fronttitle}[2]{%
    \begingroup
        \def\xmu@inabstract{true}%
        \titleformat{\chapter}{\centering#1\linespread{1}\selectfont\setlength{\baselineskip}{36bp}}{}{0pt}{}%
        \titlespacing{\chapter}{0pt}{0pt}{24bp}%
        \chapter*{#2}%
    \endgroup
    \markboth{#2}{#2}%
    \pagestyle{xmu@front}\thispagestyle{xmu@front}%
}
\renewcommand{\chaptermark}[1]{%
    \markboth{\ifcsdef{xmu@arabicchapter}{\thechapter}{第\zhnumber{\thechapter}章}\quad #1}{}%
}
\renewcommand{\chapter}{
    \clearpage
    \ifthenelse{\isundefined{\xmu@inabstract}}{
        \ifthenelse{\isundefined{\xmu@print}}{}{\cleardoublepage}
    }{}
    \xmu@bodyparagraphs
    \pagestyle{fancy}
    \fancyhf{}
    \renewcommand{\headrulewidth}{0.4pt}
    \fancyhead[OC]{\rm\zihao{-5}\leftmark}  %奇数页眉
    \fancyhead[EC]{\rm\zihao{-5}\xmu@title} %偶数页眉
    \fancyfoot[C]{\rm\zihao{-5}\thepage}    %页脚
    \global\@topnum\z@
    \@afterindenttrue
    \secdef\@chapter\@schapter
}

\fancypagestyle{plain}{\fancyhf{}\fancyhead[OC]{\rm\zihao{-5}\leftmark}\fancyhead[EC]{\rm\zihao{-5}\xmu@title}\fancyfoot[C]{\rm\zihao{-5}\thepage}\renewcommand{\headrulewidth}{0.4pt}}
